% ============================================================
% Preámbulo: paquetes y configuración global
% ============================================================

% --- Idioma y codificación ---
\usepackage[utf8]{inputenc}      % entrada UTF-8 (si usás xelatex/lualatex, quitar)
\usepackage[T1]{fontenc}         % codificación de fuente con acentos
\usepackage[spanish,es-tabla]{babel} % español (es-tabla: "Tabla" en vez de "Cuadro")

% --- Tipografía Arial (Reglamento de Titulación, Art. 22) ---
% Con pdflatex no existe Arial; se usa Helvetica (helvet), métricamente
% equivalente a Arial. Para usar Arial REAL, compilar con xelatex/lualatex
% y reemplazar este bloque por:
%   \usepackage{fontspec}
%   \setmainfont{Arial}
\usepackage{helvet}
\renewcommand{\familydefault}{\sfdefault} % fuente base sans-serif (estilo Arial)

% --- Márgenes (Reglamento de Titulación, Art. 22) ---
% izquierdo 4 cm, derecho 2 cm, superior e inferior 2,5 cm,
% encabezado y pie de página a 1,25 cm del borde.
\usepackage[letterpaper,
  left=4cm, right=2cm, top=2.5cm, bottom=2.5cm,
  headsep=0.76cm, headheight=14pt, footskip=1.25cm]{geometry}

% --- Interlineado ---
\usepackage{setspace}
\onehalfspacing                  % interlineado 1.5

% --- Sangría y espacio entre párrafos ---
\usepackage{indentfirst}         % sangra también el primer párrafo

% --- Gráficos, tablas y figuras ---
\usepackage{graphicx}
\graphicspath{{figuras/}}        % carpeta de imágenes
\usepackage{float}
\usepackage{booktabs}            % tablas profesionales
\usepackage{array}
\usepackage{longtable}
\usepackage{caption}
% Estilo APA 7: "Tabla N" / "Figura N" en negrita y, en la línea siguiente,
% el título en cursiva (Guía APA — Biblioteca U. Americana).
\captionsetup{labelfont=bf, textfont=it, labelsep=newline,
  justification=raggedright, singlelinecheck=false}

% --- Matemáticas ---
\usepackage{amsmath,amssymb}

% --- Código fuente ---
\usepackage{listings}
\usepackage{xcolor}
\lstset{
  basicstyle=\ttfamily\footnotesize,
  keywordstyle=\color{blue},
  commentstyle=\color{gray},
  stringstyle=\color{teal},
  numbers=left,
  numberstyle=\tiny,
  frame=single,
  breaklines=true,
  showstringspaces=false,
  captionpos=b
}

% --- Bibliografía estilo APA (biblatex + biber) ---
\usepackage[backend=biber,style=apa,sorting=nyt,language=spanish]{biblatex}
\DeclareLanguageMapping{spanish}{spanish-apa}
\setlength{\bibhang}{1.27cm}      % sangría francesa APA (1,27 cm = ½ pulg.)

% --- Encabezados y pies de página ---
\usepackage{fancyhdr}
\pagestyle{fancy}
\fancyhf{}
\fancyhead[L]{\nouppercase{\leftmark}}
\fancyhead[R]{\thepage}
\renewcommand{\headrulewidth}{0.4pt}

% --- Cita textual larga (>40 palabras), estilo APA ---
% Bloque con sangría de 1,27 cm a la izquierda, sin comillas ni cursiva.
% Uso:
%   \begin{citalarga} ... texto de la cita ... \parencite[p.~90]{clave} \end{citalarga}
\newenvironment{citalarga}{%
  \begin{list}{}{\setlength{\leftmargin}{1.27cm}\setlength{\rightmargin}{0cm}}%
  \item[]}{%
  \end{list}}

% --- Hipervínculos (debe ir casi al final) ---
\usepackage[hidelinks]{hyperref}
\hypersetup{
  pdftitle={Tesis de grado},
  pdfauthor={Oscar G. Pavón},
  bookmarksnumbered=true
}

% --- Numeración de capítulos en el índice ---
\setcounter{tocdepth}{2}
\setcounter{secnumdepth}{3}
